% =====================================================================
%  2bat-ud2-9.tex  —  SESSIÓ 9: Maximitzar la cobertura social (estudi de cas)
% =====================================================================
\horatitol{Sessió 9 — Maximitzar la cobertura social}%
{Estudi de cas (Activitat 3): trobar la combinació d'activitats que atén el màxim de persones i justificar per què és la millor decisió.}

\subsection{Idea clau}

Apliquem el mètode complet a un cas real. Sobre la regió factible que ja coneixem,
definim una \textbf{funció objectiu} que mesura la \textbf{cobertura social} (les
persones ateses) i busquem la combinació que la \textbf{maximitza}. El més important
no és només el càlcul: és \textbf{justificar} per què aquella és la millor decisió
possible amb els recursos disponibles.

\begin{idea}
\textbf{El cas.} L'entitat fa \textbf{tallers} ($x$) i \textbf{sortides} ($y$), amb la
regió factible coneguda (vèrtexs $(0,0)$, $(12,0)$, $(6,3)$, $(0,5)$). Cada
\textbf{taller} atén $20$ persones i cada \textbf{sortida}, $30$. La cobertura total és
\[
P=20x+30y \quad(\text{persones ateses}).
\]
\textbf{Mètode:} (1) avaluar $P$ a cada vèrtex; (2) triar el valor més gran; (3)
interpretar i \textbf{justificar} la decisió (quantes activitats, quantes persones,
quins recursos s'esgoten).
\end{idea}

\subsection{Exemples}

\begin{exemple}[avaluar la cobertura als vèrtexs]
Avaluem $P=20x+30y$ a cada vèrtex de la regió factible:
\begin{center}
\begin{tabular}{lcl}
\toprule
\textbf{Vèrtex} & $P=20x+30y$ & \\
\midrule
$(0,0)$ & $0$ & \\
$(12,0)$ & $20\cdot 12+30\cdot 0=240$ & \\
$(6,3)$ & $20\cdot 6+30\cdot 3=120+90=210$ & \\
$(0,5)$ & $20\cdot 0+30\cdot 5=150$ & \\
\bottomrule
\end{tabular}
\end{center}
El màxim és $P=240$, al vèrtex $(12,0)$.
\end{exemple}

\begin{exemple}[interpretar i justificar]
L'òptim $(12,0)$ vol dir \textbf{fer $12$ tallers i cap sortida}, que atén
\textbf{$240$ persones}. \textbf{Justificació:} amb el pressupost disponible, els
tallers «rendeixen» més cobertura per euro que les sortides (un taller és més barat i
atén prou gent), de manera que concentrar els recursos en tallers maximitza les
persones ateses. En aquesta decisió s'esgota el pressupost ($12$ tallers el consumeixen
tot). És la millor combinació \emph{possible}, perquè qualsevol altra de la regió atén
menys persones.
\end{exemple}

\subsection{Exercicis resolts}

\begin{exresolt}[maximitzar i interpretar]
Sobre la mateixa regió ($(0,0)$, $(12,0)$, $(6,3)$, $(0,5)$), suposa ara que cada
\textbf{taller} atén $15$ persones i cada \textbf{sortida}, $40$. Quina és la decisió
òptima?
\solucio
Funció objectiu $P=15x+40y$. Avaluem als vèrtexs:
\[
P(0,0)=0,\quad P(12,0)=180,\quad P(6,3)=90+120=210,\quad P(0,5)=200.
\]
El màxim és $P=210$, al vèrtex $(6,3)$: fer \textbf{$6$ tallers i $3$ sortides}, que
atén $210$ persones. Ara, com que les sortides atenen molta més gent, la millor
decisió ja no és tot tallers, sinó una combinació.
\resposta{Òptim $(6,3)$: $6$ tallers i $3$ sortides, $210$ persones ateses.}
\end{exresolt}

\begin{exresolt}[justificar el canvi d'òptim]
Per què, en canviar les persones ateses per activitat, l'òptim s'ha desplaçat de
$(12,0)$ a $(6,3)$?
\solucio
La funció objectiu mesura la cobertura, i el seu pendent depèn de quanta gent atén
cada tipus d'activitat. Quan les sortides passen a atendre molta més gent ($40$ contra
$30$), «val la pena» dedicar-hi recursos, i el vèrtex que combina tallers i sortides
$(6,3)$ supera el de només tallers. L'òptim és sensible als valors de la funció
objectiu: el mètode (avaluar als vèrtexs) és el mateix, però el guanyador canvia.
\resposta{Perquè el pes de cada variable a la funció objectiu ha canviat, i això
desplaça el vèrtex òptim.}
\end{exresolt}

\subsection{Exercicis per fer}

\noindent\textit{Regió factible del cas: vèrtexs $(0,0)$, $(12,0)$, $(6,3)$, $(0,5)$.}

\begin{exercicis}
\begin{exlist}
  \item Cada taller atén $20$ persones i cada sortida $30$ ($P=20x+30y$). Avalua $P$
        a cada vèrtex i digues quina és la decisió òptima i quantes persones atén.
  \item Ara cada taller atén $10$ persones i cada sortida $50$ ($P=10x+50y$). Troba
        l'òptim i interpreta'l.
  \item Amb $P=25x+25y$ (totes dues activitats atenen igual), avalua als vèrtexs. Què
        observes? Hi ha un únic òptim?
  \item Per a l'òptim que has trobat a l'exercici 1, digues quins recursos s'esgoten
        (pressupost $x+2y\le 12$, material $x+3y\le 15$) i quins sobren.
  \item \textbf{(Justificació)} Tria l'òptim de l'exercici 2 i redacta, en dues o tres
        frases, una conclusió justificada com si fos la proposta de l'entitat: quantes
        activitats de cada tipus, quantes persones ateses i per què és la millor opció.
  \item \textbf{(Decisió amb criteri)} Si l'entitat valorés també arribar a col·lectius
        diferents (i no només el nombre total de persones), creus que l'òptim purament
        matemàtic sempre seria la millor decisió \emph{real}? Comenta-ho breument.
\end{exlist}
\end{exercicis}

% ---------------------------------------------------------------------
\fullsolucions{Sessió 9}
\begin{solucions}
\begin{sollist}
  \item $P(0,0)=0$, $P(12,0)=240$, $P(6,3)=210$, $P(0,5)=150$. Òptim $(12,0)$: $12$
        tallers i cap sortida, $240$ persones.
  \item $P=10x+50y$: $P(0,0)=0$, $P(12,0)=120$, $P(6,3)=60+150=210$, $P(0,5)=250$.
        Òptim $(0,5)$: $5$ sortides i cap taller, $250$ persones (ara les sortides
        atenen molta més gent).
  \item $P=25x+25y$: $P(0,0)=0$, $P(12,0)=300$, $P(6,3)=225$, $P(0,5)=125$. Òptim
        $(12,0)$ amb $P=300$; com que $P=25(x+y)$, maximitzar $P$ equival a maximitzar
        $x+y$, i el vèrtex amb més activitats totals és $(12,0)$.
  \item A $(12,0)$: pressupost $12+0=12$ (s'esgota), material $12+0=12\le 15$ (en
        sobren $3$).
  \item Resposta oberta. Ha d'indicar $5$ sortides, $0$ tallers, $250$ persones, i
        justificar que amb aquests valors de cobertura les sortides són les que més
        gent atenen.
  \item Resposta oberta. Idea esperada: l'òptim matemàtic maximitza una mesura
        concreta (persones totals); si hi ha altres criteris (equitat, arribar a
        col·lectius vulnerables), la millor decisió real pot diferir, i cal afegir
        aquests criteris al model o ponderar-los.
\end{sollist}
\end{solucions}
